% Титульный лист автореферата (по шаблону МФТИ; нумерация страницы не ставится)
\thispagestyle{empty}

\begin{flushright}
{\large\bfseries\itshape На правах рукописи}\\[0.3em]
\ifnumequal{\value{showperssign}}{0}{%
    \rule[0pt]{0pt}{1.5cm}%
}{%
    \includegraphics[height=1.5cm]{personal-signature.png}%
}
\end{flushright}

\vspace{0pt plus2fill} %число перед fill = кратность относительно некоторого расстояния fill, кусками которого заполнены пустые места
\begin{center}
\textbf{\large \thesisAuthor}
\end{center}

\vspace{0pt plus2fill} %число перед fill = кратность относительно некоторого расстояния fill, кусками которого заполнены пустые места
\begin{center}
\textbf{\Large \thesisTitle}

\vspace{0pt plus2fill} %число перед fill = кратность относительно некоторого расстояния fill, кусками которого заполнены пустые места
\textbf{\large Специальность: \thesisSpecialtyNumber\ \thesisSpecialtyTitle}

\ifdefined\thesisSpecialtyTwoNumber
\textbf{\large Специальность: \thesisSpecialtyTwoNumber\ \thesisSpecialtyTwoTitle}
\fi

\vspace{0pt plus2fill} %число перед fill = кратность относительно некоторого расстояния fill, кусками которого заполнены пустые места
{\Large\bfseries АВТОРЕФЕРАТ\par}
\vspace{0pt plus1fill} %число перед fill = кратность относительно некоторого расстояния fill, кусками которого заполнены пустые места
{\large\bfseries диссертации на соискание ученой степени\par \thesisDegree\par}
\end{center}

\vspace{0pt plus4fill} %число перед fill = кратность относительно некоторого расстояния fill, кусками которого заполнены пустые места
{\centering\bfseries\thesisCity~--- \thesisYear\par}

\newpage
% оборотная сторона обложки
\thispagestyle{empty}
Работа прошла апробацию в Федеральном исследовательском центре «Информатика и управление» Российской академии наук  (ФИЦ ИУ РАН).

\vspace{0.008\paperheight plus1fill}
\noindent%
\begin{tabularx}{\textwidth}{@{}lX@{}}
    \ifdefined\supervisorTwoFio
    \textbf{Научные руководители:} & \ifdefined\supervisorDead
                              \framebox{\textbf{\supervisorFio}}\textbf{\ "--- \supervisorRegalia}
                              \else
                              \textbf{\supervisorFio\ "--- \supervisorRegalia}
                              \fi
                              \par
                              \vspace{0.013\paperheight}
                              \ifdefined\supervisorTwoDead
                              \framebox{\textbf{\supervisorTwoFio}}\textbf{\ "--- \supervisorTwoRegalia}
                              \else
                              \textbf{\supervisorTwoFio\ "--- \supervisorTwoRegalia}
                              \fi
                              \vspace{0.013\paperheight}\\
    \else
    \textbf{Научный руководитель:} & \ifdefined\supervisorDead
                              \framebox{\textbf{\supervisorFio}}\textbf{\ "--- \supervisorRegalia}
                              \else
                              \textbf{\supervisorFio\ "--- \supervisorRegalia}
                              \fi
                              \vspace{0.013\paperheight}\\
    \fi
%%     Официальные оппоненты:  &
%%     \ifnumequal{\value{showopplead}}{0}{\vspace{13\onelineskip plus1fill}}{%
%%         \textbf{\opponentOneFio,}\par
%%         \opponentOneRegalia,\par
%%         \opponentOneJobPlace,\par
%%         \opponentOneJobPost\par
%%         \vspace{0.01\paperheight}
%%         \textbf{\opponentTwoFio,}\par
%%         \opponentTwoRegalia,\par
%%         \opponentTwoJobPlace,\par
%%         \opponentTwoJobPost
%%     \ifdefined\opponentThreeFio
%%         \par
%%         \vspace{0.01\paperheight}
%%         \textbf{\opponentThreeFio,}\par
%%         \opponentThreeRegalia,\par
%%         \opponentThreeJobPlace,\par
%%         \opponentThreeJobPost
%%     \fi
%%     }%
    \vspace{0.013\paperheight} \\
    \ifdefined\leadingOrganizationTitle
    \textbf{Ведущая организация:}    &
    \ifnumequal{\value{showopplead}}{0}{\vspace{6\onelineskip plus1fill}}{%
        \leadingOrganizationTitle
    }%
    \fi
\end{tabularx}
\vspace{0.008\paperheight plus1fill}

Защита состоится \textbf{\uline{\defenseDate}}~на~заседании диссертационного совета \textbf{\defenseCouncilNumber}, созданного на базе \defenseCouncilTitle,

\textbf{по адресу:} \defenseCouncilAddress.

\vspace{0.008\paperheight plus1fill}
С диссертацией можно ознакомиться в библиотеке \synopsisLibrary.

\vspace{0.008\paperheight plus1fill}
Автореферат разослан \synopsisDate~\underline{\hspace{8em}} \thesisYear~г.

%% \noindent Телефон для справок: \defenseCouncilPhone.

\vspace{0.008\paperheight plus1fill}
\noindent%
\begin{tabularx}{\textwidth}{@{}%
>{\raggedright\arraybackslash}b{18em}@{}
>{\centering\arraybackslash}X
r
@{}}
    {\bfseries Ученый секретарь\par
    диссертационного совета,\par}
    \defenseSecretaryRegalia
    &
    \ifnumequal{\value{showsecrsign}}{0}{}{%
        \includegraphics[width=2cm]{secretary-signature.png}%
    }%
    &
    \textbf{\defenseSecretaryFio}
\end{tabularx}
